%sig-alternate.tex
%Alternate ACM SIG Proceedings document using LaTeX2e
%Author: Margot Duek
%based upon LaTeX2.09 Guidelines
%Revisions: N/A
\documentclass{sig-alternate-05-2015}
\begin{document}

\setcopyright{acmcopyright}

\title{Graph prossesing
\titlenote{(Produces the permission block, and
copyright information). For use with
SIG-ALTERNATE.CLS. Supported by ACM.}}
\subtitle{Graph prossesing by using the SNAP library
\titlenote{A full version of this paper is available as
\textit{Author's Guide to Preparing ACM SIG Proceedings Using
\LaTeX$2_\epsilon$\ and BibTeX} at
\texttt{www.acm.org/eaddress.htm}}}


\numberofauthors{1}
\author{
\alignauthor
Margot Duek\titlenote{}\\
       \affaddr{Instituto Tecnologico de Estudios Superiores de Monterrey}\\
       \affaddr{Mexico City}\\
       \email{margotduek@gmail.com}
}

\maketitle
\begin{abstract}
This paper provides a set of instructions of how to import a dataset that
includes a sample of a Facebook graph from the SNAP (Stanford Network
Analysis Platform) library which was created for analysis and graph mining.

This also includes the steps of how to export the graph to GraphML, GEXF,
GDF and JSON Graph Format with GraphSON afterwards we export the graph
with each of the formats to Gephi.

By reading these paper you will also find how to create a new document on
\LaTeX\  using the ACM format which is being used on these paper.

\end{abstract}


% The code below was generated by the tool at
% http://dl.acm.org/ccs.cfm
\begin{CCSXML}
  <ccs2012>
  <concept>
  <concept_id>10011007</concept_id>
  <concept_desc>Software and its engineering</concept_desc>
  <concept_significance>500</concept_significance>
  </concept>
  </ccs2012>
\end{CCSXML}

\ccsdesc[500]{Computer systems organization~Embedded systems}
\ccsdesc[300]{Computer systems organization~Redundancy}
\ccsdesc{Computer systems organization~Robotics}
\ccsdesc[100]{Networks~Network reliability}

% End generated code


\keywords{ACM proceedings; \LaTeX; text tagging}

\section{Introduction}
A graph is a data structure with a finite set of vertex or nodes
that get together by edges, there are directed graphs, that are
the graphs wich edges point to a direction, and undirected graphs,
which edges point in both ways.

Stanford Network Analysis Platform is a library being actively developed
since 2004, this is a very efficient library wich can analyze hundreds
of millions of nodes and billions of edges, this can scale to massive networks
and is written in c++.


\section{Graphs}
We have already talked about what are graphs, what is SNAP and what can
we do when them, but we haven't talked about what can we do with
these tools or how to use them.
  
Is important to say that we will be working with the facebook graph sample,
but you can do it with any graph, just remember to change the name of the
file every time we need to make a route.

\subsection{Importing the Facebook graph}
The first thing we are going to talk about is how to import a dataset that
includes a sample of a Facebook graph. You need to get into
\texttt{https://snap.stanford.edu/data/index.html} and choose which
graph you want to use, in these case, we will get into social networks directory
and inside there we will choose Facebook, remember you can choose anything.

You will need to download the facebook\_combined.txt.gz, unzip it in the folder

where you want to use it (Don't forget to add the SNAP library in the
directory where you will work) and when you see the file, you will see a
bunch of numbers, those numbers are the graph.

Crete your main.cpp, there you will have to import the graph and check if it was
imported successfully, inside the main, you will need to put a line that should
look like these,
\texttt{PUNGraph g=TSnap::LoadEdgeList<PUNGraph>("../facebook
\_combined.txt",0,1);}
where you are declaring a new PUNGraph variable
and giving to that variable the Facebook graph, now we have imported the graph.

The only way to know if something really worked is by checking it, we
will print the amount of nodes that we have on the graph by printing
the name of your graph.GetNodes(), in this case, g->GetNodes(); and it
it prints the amount of nodes that the SNAP's website says that means you
have successfully imported the graph.


\subsection{Exporting the graph}
Now is time to talk about how to export the graph, we will
export the graph in four different ways, GraphML, GEXF,
GDF and JSON Graph Format with GraphSON, each of them have their own
implementation but these time we will talk about the general
process we should make and after that, you should be able
to make the same thing with the other ones.

Basically what we need to do, is to make a new ofstream file
and name it with the name you like and we will start printing
what we want inside that file(in each case we have
a different file to export the graph in each format),
this method should have a condition or a cycle so
that all the nodes can be printed  of that file.

GraphML is the less efficient format of exporting from
the four we tasted, it took 26 milliseconds with a
linear temporal complexity in the worst of the cases
because it have two for's inside of an if.

Executing GEXF function took us a total of 22 seconds,
this also have in the worst of the cases a temporal complexity
that is linear because inside of an if we have two separated
for's to convert each node.

The total time of converting the facebook sample graph
to GDF format was 12 milliseconds, which makes these
function the one with the best time. These function
have the same temporal complexity in the worst of the cases
as the previous ones because is almost the same thing
to do.

The last function implemented was to export the graph
to a JSON Graph Format with GraphSON which also has
a very good time, we have got 16 seconds. The
implementation of these function is based on the same
architecture than the other ones, an if condition
and two for cycles to export each node.



\subsection{LaTex}
A very big problem for people who is not used to designing and thinking
of how should things go on an investigation paper is that we have to
think where to put each of the things in the document, like the date,
the name, personal information and other things.

For all the people that is most of all interested in the real
content, you can use \LaTeX\ , this is a tool that allows people
to just write content and don't worry about designing problems,
and let the designers take care of that.

The only thing we should do is (if you are in Linux, if you are not,
you should find a paper with a tutorial for your operating system)
install any tool for \LaTeX\ , in my case I used Texlive, after
doing that, you should open your favorite text editor and start
writing your paper, in my case I used a template from ACM
and I edited it.

To see what you have created, you need to make your .tex file into a
.pdf by doing ctrl + Alt + B, if when you do that nothing happens,
that means you have a compilation error, you should check the syntaxis
of everything that has been written in the document.


\subsection{Gephi}
Gephi is an interactive visualization and exploration platform
for all kinds of networks and complex systems,
dynamic and hierarchical graphs.

\subsection{The code}
You can find the code in https://github.com/margotduek/
AnalisisDisenioDeAlgoritmos/tree/master/tarea4/prueba
which is a directory on a GitHub repository.

\begin{figure}
\centering
\includegraphics{gephi}
\caption{A sample Gephi graph}
\end{figure}

\section{Conclusions}
By using any library, in this case, we used SNAP but there
are a lot of libraries that we can use for any subject or
a project, we can make bigger and best applications in less time
and in some cases, using a library could be more efficient
because of the way the algorithms are programmed.


%ACKNOWLEDGMENTS are optional
\section{Acknowledgments}
In this investigation and while using the graph library
I learned a lot of things, especially I learned how to use
\LaTeX\ for a paper with the ACM standard or for anything that I
need to make an investigation paper ir report.

I think this is a very useful tool for all my career and all my life
because for me, all my life it has been very difficult to find out where
to put each information, and by using \LaTeX\ is easier and
faster for me to make my reports.

I also learned how easy is to use a library, I think in most of the
cases, everyting you need is time because you need to read a lot
of things but if you read all the documentation you can
do almost everything.

%
% The following two commands are all you need in the
% initial runs of your .tex file to
% produce the bibliography for the citations in your paper.
\bibliographystyle{abbrv}
\bibliography{sigproc}  % sigproc.bib is the name of the Bibliography in this case
% You must have a proper ".bib" file
%  and remember to run:
% latex bibtex latex latex
% to resolve all references
%
% ACM needs 'a single self-contained file'!
%
%APPENDICES are optional
%\balancecolumns
\appendix
%Appendix A
\section{Headings in Appendices}
\subsection{Introduction}
\subsection{Graphs}
\subsubsection{Importing the Facebook graph}
\subsubsection{Exporting the graph}
\subsubsection{\LaTeX;}
\subsubsection{Gephi}
\subsubsection{The code}
\subsection{Conclusions}
\subsection{Acknowledgments}
%\balancecolumns % GM June 2007
% That's all folks!
\end{document}
